% ============================================================================================================
%  Exposé zur Abschlussarbeit im FG Partizipative IT-Gestaltung  –  DEUTSCH
%  Version 1.0 (2026)
%  Schriftart: Atkinson Hyperlegible | APA-Literaturverzeichnis | A4
% ============================================================================================================
\documentclass[12pt, a4paper]{article}

% ── Kodierung & Sprache ─────────────────────────────────────
\usepackage[utf8]{inputenc}
\usepackage[T1]{fontenc}
\usepackage[ngerman]{babel}
\usepackage{csquotes}
\usepackage{float}

% ── Font: Atkinson Hyperlegible ───────────────────────────────────
\usepackage{atkinson}
\renewcommand{\familydefault}{\sfdefault}

% ── Mikrotypografie ─────────────────────────────────────────
\usepackage[activate={true,nocompatibility},
            final, tracking=true, kerning=true,
            spacing=true, factor=1100, stretch=10, shrink=10]{microtype}

% ── Seitenlayout ────────────────────────────────────────────
\usepackage[
  a4paper,
  top=2.5cm, bottom=2.5cm,
  left=2.5cm, right=2.5cm
]{geometry}
\usepackage{setspace}
\onehalfspacing

% ── Kopf- & Fußzeilen ───────────────────────────────────────
\usepackage{fancyhdr}
\pagestyle{fancy}
\fancyhf{}
\fancyhead[L]{\small\textit{Exposé – Abschlussarbeit FG PIT}}
\fancyhead[R]{\small\textit{\dieAutorin}}
\fancyfoot[C]{\thepage}
\renewcommand{\headrulewidth}{0.4pt}

% ── Grafik & Farbe ──────────────────────────────────────────
\usepackage{graphicx}
\usepackage[table]{xcolor}
\definecolor{viridis-purple}{RGB}{68, 55, 129}

% ── Abschnittsformatierung ──────────────────────────────────
\usepackage{titlesec}
\titleformat{\section}
  {\large\bfseries\color{viridis-purple}}{\thesection}{1em}{}[\titlerule]
\titleformat{\subsection}
  {\normalsize\bfseries}{\thesubsection}{1em}{}

% ── Tabellen ────────────────────────────────────────────────
\usepackage{booktabs}
\usepackage{tabularx}
\usepackage{array}

% ── Gantt / Zeitplan ────────────────────────────────────────
\usepackage{pgfgantt}

% ── Hyperlinks ──────────────────────────────────────────────
\usepackage[
  colorlinks=true,
  linkcolor=viridis-purple,
  citecolor=viridis-purple,
  urlcolor=viridis-purple
]{hyperref}

% ── Literaturverzeichnis (APA 7. Auflage) ───────────────────
% Sie können alternativ auch mit den Zitationsstilen der IEEE oder 
% nach der Chicago Zitierweise arbeiten; dann bitte hier entsprechend
% ändern
\usepackage[
  backend=biber,
  style=apa,
  sortcites=true,
  sorting=nyt
]{biblatex}
\addbibresource{references.bib}

% ── Hilfspakete ─────────────────────────────────────────────
\usepackage{enumitem}
\usepackage{parskip}
\usepackage{lipsum}

% ── Eigene Befehle ──────────────────────────────────────────
\newcommand{\dieAutorin}{Linnea Herzer}          % Einmal ändern, überall sichtbar
\newcommand{\derTitel}{Analyse der User Experience eines textilen Interfaces zur Musiksteuerung}
\newcommand{\dieBetreuerIn}{Prof. Dr. Claude Draude}
\newcommand{\dasInstitut}{Fachgebiet Partizipative IT-Gestaltung}
\newcommand{\dieUniversitaet}{Fachbereich 16 - Elektrotechnik/ Informatik}
\newcommand{\dasDatum}{\today}

% ============================================================
\begin{document}

% ── Titelseite ──────────────────────────────────────────────
\begin{titlepage}
  \centering

  \vspace*{1cm}
  \includegraphics[width=10cm]{uni-logo.png}   % <-- Logo-Datei einfügen
  \par\vspace{0.5cm}

  {\large \dieUniversitaet \par}
  {\normalsize \dasInstitut \par}
  \vspace{1.5cm}

  \rule{\linewidth}{0.5pt}\\[0.4cm]
  {\LARGE \bfseries \derTitel \par}
  \rule{\linewidth}{0.5pt}\\[1cm]

  {\large \textbf{Exposé} \par} % korrektes auswählen
  {\normalsize Bachelor of Science \par} % korrektes auswählen

  \vfill

  \begin{tabular}{rl}
    \textbf{Verfasst von:}   & \dieAutorin \\[4pt]
    \textbf{Betreuung:}      & \dieBetreuerIn \\[4pt]
    \textbf{Datum:}          & \dasDatum \\
  \end{tabular}

  \vspace{1cm}
\end{titlepage}

% ── Inhaltsverzeichnis ──────────────────────────────────────
\tableofcontents
\newpage

% ====================================================================================================================
%  Generelle Hinweise zum wissenschaftlichen Arbeiten finden Sie in folgendem Moodle-Kurs: 
%  https://moodle.uni-kassel.de/course/view.php?id=17394 (Wissenschaftliches Arbeiten 101 - Fachbereich 16)
% ====================================================================================================================

% ============================================================
%  1.  PROBLEMSTELLUNG
% ============================================================
\section{Problemstellung \& Motivation}

\subsection{Problemstellung}
Smart Textiles werden zunehmend in verschiedenen Bereichen wie Gesundheitswesen, Bildung und Mode eingesetzt. Über die Jahre haben sich Smart Textiles dahin entwickelt, dass die Sensoren und andere Technik kaum sichtbar sind und somit das Textil selbst als Interface für den Nutzer dient \parencite{ruckdashel2022smart}. Solche Interaktionen lösen Gefühle bei den Nutzenden eines Produkts aus, weshalb es wichtig ist die User Experience zu untersuchen um mehr Menschen zu erreichen. Ein Problem in der Forschung ist, dass bisher größtenteils die pragmatische Dimension untersucht wurde und weniger die hedonische Dimension \parencite{guennes2025smart}. Daher wird in dieser Arbeit die hedonische Dimension eines selbst entwickelten Smart Textiles zur Musiksteuerung analysiert. Dies geschieht anhand einer vorhandenen UX Bewertungsmethode, den AttrakDiff Fragebogen.

\subsection{Motivation}
Ich habe dieses Thema gewählt, weil es zu meinen Hobbys passt und kreative und technische Aspekte verbindet. Vor allem bin ich gespannt, wie ein selbst gehäkelter Prototyp bei Nutzern ankommt. Außerdem habe ich bereits in anderen Modulen des Fachgebiets PIT bemerkt, dass ich mich für Human-Computer-Interaction und User Experience interessiere. Im Bereich Elektrotechnik habe ich bisher noch nicht viele praktische Projekte gemacht, deswegen freue ich mich darauf etwas mit dem Arduino Lilypad zu bauen.


% ============================================================
%  2.  VERWANDTE ARBEITEN
% ============================================================
\section{Verwandte Arbeiten}

\subsection{Thema 1: E-Textiles}

E-Textiles sind Textilien, in denen Elektronik wie Sensoren und leitfähiges Garn integriert sind. Sie sind eine Untergruppe der Smart Textiles, die Reize der Umwelt wahrnehmen und darauf reagieren. Das Arduino LilyPad, entwickelt von Leah Buechley, ist ein Mikrocontroller speziell für E-Textiles \parencite{buechley2008lilypad}. Es wurden bereits einige E-Textiles in vielen verschiedenen Bereichen entwickelt \parencite{younes2023smart}. Im Bereich Musik gibt es zum Beispiel ein gestricktes Keyboard mit Sensoren, das bei Berührung oder Annäherung Töne erzeugt \parencite{wicaksono2020knittedkeyboard}.

\subsection{Thema 2: User Experience Evaluation}

Nach der International Organization for Standardization beschreibt User Experience die Wahrnehmungen und Reaktionen einer Person, die aus der tatsächlichen oder erwarteten Nutzung eines Produkts, Systems oder einer Dienstleistung entstehen \parencite{iso9241-210}. Hassenzahls UX-Modell zur Beschreibung der User Experience unterteilt die User Experience in die pragmatische Qualität und die hedonische Qualität. Bei der pragmatischen Qualität wird vor allem die Funktionalität untersucht und die hedonische Qualität geht mehr auf die subjektive Wahrnehmung des Produkts ein \parencite{hassenzahl2007hedonic}. Es existieren bereits standardisierte Methoden zur Auswertung der User Experience, zum Beispiel die Fragebögen UEQ und AttrakDiff. Diese Fragebögen wurden auch bereits in verschiedenen Bereichen verwendet \parencite{diaz2019standardized}. Bei Smart Wearables handelt es sich um eine spezielle Form der Mensch-Computer-Interaktion, weil dort der Körper ein Teil des Interfaces wird. Daher ist die Auswertung der User Experience bei Smart Wearables besonders wichtig \parencite{guennes2025smart}.

\subsection{Forschungslücke}

Obwohl es bereits E-Textiles in verschiedenen Bereichen gibt und auch Methoden zur User Experience Auswertung existieren, wurde die User Experience von Smart Textiles wenig untersucht. Vor allem die hedonische Qualität wurde dabei seltener berücksichtigt. In den meisten Forschungen zu Smart Textiles geht es nur um die technische Umsetzung \parencite{guennes2025smart}. Daher wird in dieser Bachelorarbeit der Prototyp eines E-Textiles untersucht, um Erkenntnisse über dessen User Experience zu gewinnen. Die Ergebnisse beziehen sich ausschließlich auf diesen Prototyp und können nicht auf E-Textiles im Allgemeinen übertragen werden.

% ============================================================
%  3.  FORSCHUNGSFRAGE
% ============================================================
\section{Forschungsfrage}

Aus der Problemstellung ergibt sich die folgende Forschungsfrage:
\begin{quote}
  \textbf{FF:} \textit{Wie bewerten Nutzer die hedonische Dimension eines Smart Textiles zur Musiksteuerung?}
\end{quote}

Um die Forschungsfrage zu beantworten werden zwei verschiedene Perspektiven betrachtet: Welche Aspekte werden besonders positiv und welche besonders negativ bewertet?

% ============================================================
%  4.  METHODIK
% ============================================================
\section{Methodik}

\subsection{Forschungsansatz und geplante Vorgehensweise}

Es handelt sich um einen empirischen Ansatz, da Nutzerstudien mit dem Prototyp durchgeführt werden. Um die Forschungsfrage zu beantworten, werden Daten benötigt um herauszufinden wie die Nutzer das textile Interface zur Musiksteuerung bewerten. Dafür wird in der Studie ein Usability Test durchgeführt und anschließend ein standardisierter Fragebogen ausgefüllt.

\subsection{Studiendesign} 

\begin{description}[style=nextline, leftmargin=2em]
  \item[Teilnehmende]
  Die Zielpopulation der Studie sind Erwachsene ab 18 Jahren. Die Teilnehmenden werden überwiegend aus dem persönlichen Umfeld rekrutiert. Außerdem werden noch einzelne Personen außerhalb des persönlichen Umfelds rekrutiert, um das Risiko sozialer Erwünschtheit bei der hedonischen Qualität zu reduzieren. Einschlusskriterien sind ein Mindestalter von 18 Jahren und die regelmäßige Nutzung digitaler Musiksteuerung. Personen unter 18 Jahren werden von der Studie ausgeschlossen. Da der Umfang einer Bachelorarbeit begrenzt ist, ist ein Stichprobenumfang von zehn Personen geplant.

  \item[Materialien \& Apparatus]
  Für die Nutzerstudie wird ein gehäkeltes Armband zur Musiksteuerung als Prototyp gebaut. Dafür wird aus Baumwolle ein Armband gehäkelt und anschließend ein geeigneter Mikrocontroller darin vernäht. Leitfähiges Garn dient zur Herstellung der elektrischen Verbindungen. Im weiteren Verlauf wird noch geprüft, ob ein Mikrocontroller mit integrierter Bluetooth Funktion oder ein extra Bluetooth-Modul verwendet wird. Bei einem separaten Bluetooth-Modul müsste entsprechend auf die Größe geachtet werden und ob es mit leitfähigem Garn angenäht werden kann. Zur Programmierung des Mikrocontrollers wird die Arduino IDE verwendet.
  In der Nutzerstudie kommt der Fragebogen AttrakDiff zur Auswertung der User Experience zum Einsatz.

  \item[Durchführung]
   Begrüßung \& Briefing → Einwilligungserklärung → Bearbeitung der Aufgaben → Beobachtung und Dokumentation → Ausfüllen des Fragebogens → Debriefing

  \item[Erhobene Variablen]
  Variablen: Pragmatische Qualität, Hedonische Qualität, Attraktivität\\
  Operationalisierung: Verwendung des AttrakDiff Fragebogens

  \item[Auswertung]
  Die Ergebnisse des Fragebogens werden mit geeigneten Diagrammen visualisiert. Der Fragebogen wird deskriptiv ausgewertet, durch Berechnung des Mittelwerts und der Standardabweichung. Dabei liegt der Fokus auf den Skalen Hedonische Qualität-Stimulation (HQ-S) und Hedonische Qualität-Identität (HQ-I). Die Beobachtungen während der Nutzung werden ergänzend qualitativ ausgewertet, um Auffälligkeiten zu erfassen und die Ergebnisse des Fragebogens besser einzuordnen.
    
\end{description}

\subsection{Ethische Überlegungen}

Die Teilnahme an der Studie erfolgt freiwillig und kann jederzeit ohne Angabe von Gründen abgebrochen werden. Während der Studie können personenbezogene Daten erhoben werden. Daher erhält jeder Proband vor der Studie eine Einwilligungserklärung, in der über die erhobenen Daten und deren Verarbeitung informiert wird. Dafür wird die bereitgestellte Einwilligungserklärung von Moodle verwendet und entsprechend an die Studie angepasst. Die Studie wird nur nach Einwilligung des Probanden durchgeführt. Der Fragebogen und die Beobachtungsnotizen werden mit einer Probanden-ID versehen, um sie gemeinsam auswerten zu können.

% ============================================================
%  5.  GEPLANTE GLIEDERUNG
% ============================================================
\section{Geplante Gliederung (+ Seitenanzahl)}

\begin{enumerate}
    \item Einleitung (2)
    \item Theoretische Grundlagen (12)
    \begin{enumerate}
        \item User Experience
        \begin{enumerate}
            \item Definition von UX
            \item Modelle der User Experience
            \item Methoden zur UX-Messung
        \end{enumerate}
        \item Tangible Interaction und Materialität
        \begin{enumerate}
            \item Grundlagen der Tangible Interaction
            \item Horneckers Framework
        \end{enumerate}
        \item Textile Interfaces
        \begin{enumerate}
            \item Grundlagen und Technologien
            \item eTextiles und LilyPad
            \item Interaktionsformen
        \end{enumerate}
        \item Musiksteuerung
        \begin{enumerate}
            \item Grundlagen und Anforderungen
            \item Interaktionsformen
        \end{enumerate}
        \item Stand der Forschung
    \end{enumerate}
    \item Praktischer Teil (16)
    \begin{enumerate}
        \item Methodik
        \item Versuchsaufbau
        \item Ergebnisse
        \item Diskussion
    \end{enumerate}
    \item Fazit (2)
\end{enumerate}


% ============================================================
%  6.  ARBEITSPLAN & ZEITPLAN
% ============================================================
\section{Arbeitsplan \& Zeitplan}

\medskip

% ── Gantt Chart ─────────────────────────────────────────────

    
    
    % use landscape format iff chart becomes too wide
    %\begin{landscape}
    
    \begin{figure}[H]
    \begin{ganttchart}[
      hgrid, vgrid,
      x unit=0.7cm,
      y unit chart=0.5cm,
      newline shortcut=true,
      group label node/.append style={align=left}
    ]{1}{9}
      % Titelzeile mit Wochen
      \gantttitle{W1}{1}
      \gantttitle{W2}{1}
      \gantttitle{W3}{1}
      \gantttitle{W4}{1}
      \gantttitle{W5}{1}
      \gantttitle{W6}{1}
      \gantttitle{W7}{1}
      \gantttitle{W8}{1}
      \gantttitle{W9}{1}\\
    
      % Aufgabenbalken
      \ganttbar{Finalisierung des Prototyps und technische Tests}{1}{2} \\
      \ganttbar{Vorbereitung der Nutzerstudie}{1}{2} \\
      \ganttbar{Durchführung der Nutzerstudie}{3}{4} \\
      \ganttbar{Auswertung der AttrakDiff-Daten}{4}{5} \\
      \ganttbar{Schreiben der Bachelorarbeit}{1}{8} \\
      \ganttmilestone{Erste vollständige Version}{7} \\
      \ganttbar{Korrekturlesen und Überarbeiten}{8}{9}
    \end{ganttchart}
    \end{figure}
    
    %\end{landscape}
    
\medskip

\textbf{Meilensteine:}
\begin{itemize}[leftmargin=2em]
  \item \textbf{Vor Anmeldung:} Abschluss der Literaturrecherche, Planung und Entwicklung des Prototyps
  \item \textbf{Woche 2:} Technische Tests und Vorbereitung der Nutzerstudie
  \item \textbf{Woche 4:} Durchführung der Nutzerstudie mit anschließender Datenerhebung.
  \item \textbf{Woche 5:} Auswertung der Daten
  \item \textbf{Woche 7:} Erste vollständige Version der Bachelorarbeit
  \item \textbf{Woche 9:} Überarbeitung der Bachelorarbeit und Abschluss der finalen Version.
\end{itemize}

% ============================================================
%  LITERATURVERZEICHNIS
% ============================================================
\newpage
\printbibliography[heading=bibintoc, title={Literaturverzeichnis}]

% ============================================================
%  ANHANG
% ============================================================
\appendix
\section{Anhang}

\subsection{Online-Fragebogen zur UX Auswertung}
Der AttrakDiff-Fragebogen ist unter folgendem Link erreichbar:
\href{https://attrakdiff.de/}{AttrakDiff Online-Fragebogen}

\end{document}
% ============================================================
%  KOMPILIERUNG:  latexmk -pdf -interaction=nonstopmode thesis_proposal_hci_DE.tex
% ============================================================